% 第5章 总结与展望

\section{工作总结}

本文围绕"把游戏代码驱动的数据合成方法迁移到真实 3D 场景"这条主线，在方法迁移、管线构建、推理优化和评估框架四个层面开展了系统性工作。

在方法迁移层面，本文首先验证了 Code2Logic 在 2D 游戏领域的有效性（30 款游戏、1397 条 QA），然后将这套模板化数据合成的思路搬到了三维场景理解任务中。通过 Concept-Graphs 管线，在 Replica 数据集的 7 个室内场景上实现了从 RGB-D 图像序列到结构化 QA 的自动标注。进一步地，为了给推断标注提供可量化的精度基准，搭建了基于 Isaac Sim 的合成真值验证系统，生成了 20 个场景、795 条带有精确 ground truth 的 QA 数据。三层递进——方法验证、真实迁移、真值校准——三类数据累计 2395 条 QA，覆盖五种推理类型。需要指出的是，Code2Logic 的 1397 条 QA 直接使用了 Tong 等人公开发布的 GameQA 数据集，本文的创造性贡献在于自行构建了 Concept-Graphs 管线的 370 条 QA 和 Isaac Sim 的 795 条 QA。

在管线构建层面，完成了一条从 RGB-D 图像序列和相机轨迹到结构化 QA 数据集的六步端到端流水线。输入经过 GradSLAM 三维重建、Grounded-SAM 二维分割、CF-SLAM 三维建图后得到物体级别的场景表示，再由 LLaVA 为每个物体生成自然语言描述，最后由 llama.cpp 在 GPU 上完成描述优化和空间关系推理。管线的关键改进集中在第五步的 prompt 设计：通过强制输出物体类别字段、设置禁用词汇表和引入容错处理，将物体类别识别率从接近零提升至约 85\%。消融实验量化了三个 prompt 组件的独立贡献，其中类别强制输出贡献最大（从 0\% 到 75\%），禁用词汇表额外贡献了约 10 个百分点。同场景新旧 prompt 的对比进一步验证了这一改进的实质效果：虽然新 prompt 的 QA 数量约降到了旧 prompt 的 30\%--35\%，但语义内容发生了从"场景氛围描述"到"物体级标注"的质的改变。

在推理优化层面，排查并解决了 Ollama 内置 CUDA 12.8 库与服务器驱动 565.57 的版本不兼容问题。通过源码编译 llama-cpp-python（链接系统 CUDA 12.6 库），将 qwen2.5 模型的全部层加载至 GPU 显存，推理吞吐量从 3 tok/s 提升至 99 tok/s（33 倍），单场景 LLM 处理时间从 15 分钟降至 50 秒。在 5 款国产开源模型上的横向对比中，Qwen2.5-3B 在速度、质量和格式合法性三个维度上综合最优。一个意外但重要的发现是，14B 模型的实际产出（8 条描述、0 条关系）远低于 3B（21 条描述、4 条关系），说明在结构化输出任务上，更大的模型安全对齐更强，反而导致更多的输出抑制。

在评估框架层面，建立了从数据量、推理类型、关系类型、标注精度和管线效率五个维度系统对比三类数据源的统一评估体系，同时也为其他数据合成方法的比较提供了一个可复用的分析模板。

\section{局限性}

这项工作存在几项需要诚实面对的局限。

最突出的瓶颈来自 LLaVA 的输入质量。7 个场景 539 条 caption 的过滤统计显示，62\% 的原始描述在进入 LLM 优化前就被预处理筛掉了——截断（以半截词开头）占 39\%，氛围词描述占 13\%。无论下游的 prompt 和推理做到什么程度，标注产量的硬上限仅有输入量的约 38\%。这个瓶颈的根源在视觉语言模型的物体裁剪阶段：当裁剪区域太小或背景噪声过多时，LLaVA 无法生成有意义的物体描述。

此外，场景多样性有限制了泛化性评估的充分性。当前的 7 个 Replica 室内场景（两个起居室、五个办公室）和 Isaac Sim 的 4 种房型模板全部属于室内静态环境，不涉及室外、工业或含有动态物体的场景。

还有两个需要后续补全的验证环节。第一，本文的实验集中于标注数据的产出过程，没有将 2395 条 QA 用于下游 VLM 的训练并评估实际效果。标注数据对模型性能的贡献是检验数据质量的终极标准，这个实验闭环在本研究中尚未完成。第二，Isaac Sim 的渲染和 ground truth 导出验证通过了，CG 管线的各步骤也在 7 个场景上独立验证了稳定性，但两个系统之间的完整自动串联流程还没有跑通。

\section{未来方向}

针对上述局限，后续工作有几个明确的方向。

在输入质量方面，基于深度信息的自适应裁剪可以在物体远近不同时动态调整裁剪框大小，兼顾近距离物体的上下文保留和远距离物体的噪声抑制。多视角描述融合也能有效滤除单帧噪声——同一物体在多个观测帧中生成的独立描述，经过交叉验证后的一致性筛选可以大幅降低幻觉率。An 等人\cite{multimodal_3d} 提出的深度与 RGB 特征联合编码方法是可直接借鉴的技术路线。

在下游验证方面，将 2395 条 QA 用于微调 Qwen2.5-VL 等视觉语言模型，在 MathVista、MMMU、GQA 等标准基准上评估微调前后的性能变化，是检验数据质量最直接的方式。Code2Logic 原始实验中使用的 GRPO 强化学习策略\cite{code2logic}——分别降低了 19.2\% 的感知错误率和 11.5\% 的推理错误率——可以引入本文的数据训练流程。

在场景扩展方面，ScanNet 和 HM3D 等更大规模的数据集是真实场景覆盖面的直接补充。Isaac Sim 中可以构建更多样化的房型模板，引入动态物体和时序连续渲染，评估管线在变化环境下的性能。Song 等人\cite{llm_planner} 关于具身代理在动态环境中多步规划挑战的研究，为动态场景的评估范式提供了参考。

真值验证链路的全自动化同样值得继续推进。将 Isaac Sim 渲染、CG 管线标注和真值对比三个环节串联后，可以在物体检测数量、类别准确率、三维位置误差和关系准确率四个指标上直接给出精确的量化结果，使"标注精度可验证"这一主张从概念落地为可操作的评估流程。
